We present a browser-based generative-design and discovery platform for atomically resolved, graphene-derived two-dimensional carbon architectures and use it for an autonomous, hypothesis-driven computational campaign on how architecture controls strength and fracture. All mechanics come from the published screened second-generation REBO potential; a PyTorch port that runs on Apple-silicon GPUs (MPS), CUDA or CPU was validated against the Atomistica Fortran reference (energies to $6\times10^{-14}$\,eV/atom and forces to $4\times10^{-13}$\,eV/\AA{} in float64; the float32 GPU mode within $7\times10^{-6}$\,eV/atom, $1.5\times10^{-4}$\,eV/\AA{} and $4\times10^{-5}$\,N/m of the reference, including bond rupture and complete stress--strain curves; 20 tests, all passed) and its neighbour list against brute-force enumeration. Five reference images (leaf venation, closed cells with an inner network, a disordered fibre net, rectilinear struts, nested rings around a void) were interpreted as design principles (hierarchy, load-path alignment, disorder, cellular partitioning, flaw halos) and translated into a parametric design space of periodic sheets. A campaign of \nRuns{} athermal quasi-static tensile simulations (2D stress in N/m) covered baselines, a reconnaissance with hierarchy as a controlled axis, five competing hypotheses tested by matched-pair discriminating experiments, deep mechanism studies, seed replicates and twelve holdout designs whose outcomes were predicted before simulation. In the model, the strength of a porous sheet is the minimum load-bearing solid fraction across the load times a net-section strength that depends on ligament width, alignment and lattice orientation ($\approx$30--32\,N/m for straight load-parallel zigzag ligaments wider than 20\,\AA, 17--19\,N/m for 10--13\,\AA{} round-pore ligaments), not on pore shape, porosity per se or hierarchy depth. Load-path orientation is the largest lever (25 vs 4\,N/m for slits parallel vs perpendicular to the load at equal porosity), and it is lost abruptly once the tips of adjacent slit rows overlap (en-echelon coalescence at 20 degrees). Hierarchy pays off only when its levels are aligned with the load; the best nested design combines load-elongated fine pores with load-parallel veins (19.4\,N/m, 2.0\,J/m$^2$ work to failure). The fracture mode is a fibre-bundle load-transfer property, flaws above $\approx$20\,\AA{} sit on a lattice-trapping plateau (18--19.5\,N/m up to 50\,\AA), and disorder, gradients and flaw halos are costs. The holdout test gave a median strength error of 14\% and the correct single-avalanche/multi-event classification in 10 of 12 cases; its misses exposed the chirality dependence of ligament strength, the additivity of alignment effects across scales and the echelon mechanism. The platform, the validated force engine, the experiment database, structures, trajectories, figures, movies and this report are delivered as one reproducible archive. All statements concern the REBO2 model under quasi-static loading; no claims about synthesizability or experimental fracture toughness are made.
